% Ch11 self-study -- "I do / You do" ladder for Zumdahl 11.1-11.3 only.
% Ten ladders, 40 questions, thirteen figures.
%
% SCOPE: the teacher ruled 11.4-11.8 out (2026-08-23), which matches
% TEXTBOOK-MAP's skip list -- colligative properties (11.4-11.7) and
% colloids (11.8) both left the CED in the 2019 redesign. What remains,
% 11.1-11.3, feeds CED 3.7 (Solutions and Mixtures) and 3.10
% (Solubility), pp. 545-563. Unlike ch 9, there is no enrichment coda
% here: the skipped material is genuinely skipped.
\documentclass{apchem-worksheet}

\apcunitword{Chapter}
\apcunit{11}
\apcunittitle{Properties of Solutions}
\apcdoctitle{Self-Study \textbullet\ Chapter 11, I~do / You~do}
\apcsubtitle{Zumdahl \S11.1--11.3 \textbullet\ PDF pp.\ 545--563}

\begin{document}
\apctitleblock

\begin{apcnote}[How to use these notes --- read this first]
Each skill is a \textbf{ladder}: a fully worked example in the
solid-framed box, then \textbf{four} YOUR TURN questions in the dashed
box --- same skill, new numbers, no help. Work all four before checking
the gray \emph{check:} line; a tracker tick needs all four right first
time.

\textbf{This document covers \S11.1--11.3 and stops there.}
Zumdahl's \S11.4--11.7 --- vapour pressure lowering, boiling point
elevation, freezing point depression, osmotic pressure, the whole of
colligative properties --- and \S11.8 on colloids are \textbf{not on the
AP CED} and are not treated here at all. That is not an oversight; it
is most of a chapter you do not have to learn.

\textbf{What is left is not small.} \S11.1--11.3 feeds CED topic
\textbf{3.7} (Solutions and Mixtures) and \textbf{3.10} (Solubility),
both inside Unit 3 --- the heaviest unit on the exam at 18--22\%.
Concentration calculations in particular turn up everywhere: in
titrations, in equilibrium, in kinetics, in electrochemistry.

\textbf{The chapter in one sentence.} A solution is a mixture at the
particle level, and every question about one is either \emph{how much
is in there} (\S11.1) or \emph{why did it dissolve at all}
(\S11.2--11.3).
\end{apcnote}

% =====================================================================
\section*{\sffamily Ladder 1 \textbullet\ Molarity}
\zsec{11.1}

Molarity is moles of solute per litre of \textbf{solution} --- not per
litre of solvent. The distinction is the whole of Ladder 1.

\[ M = \frac{\text{moles of solute}}{\text{litres of \textbf{solution}}} \]

\begin{center}
\begin{tikzpicture}[font=\sffamily\footnotesize]
  % RIGHT WAY
  \node[font=\sffamily\small\bfseries] at (2.0,4.5) {how you actually do it};
  % Volumetric flask: round bulb + narrow neck.  Drawn as a circle with
  % a white band punched through the top to erase the arc where the
  % neck joins, leaving the shoulders.  (The previous arc swept
  % 150->30 degrees, which opened upward and produced an hourglass.)
  \fill[apcgray!30] (2.0,1.30) circle (0.70);
  \draw[very thick] (2.0,1.30) circle (0.70);
  \fill[white] (1.83,1.72) rectangle (2.17,3.25);
  \fill[apcgray!30] (1.83,1.72) rectangle (2.17,2.70);
  \draw[very thick] (1.83,1.88) -- (1.83,3.25);
  \draw[very thick] (2.17,1.88) -- (2.17,3.25);
  \draw[very thick] (1.70,2.70) -- (2.30,2.70);
  \node[anchor=west,font=\tiny] at (2.38,2.70) {mark};
  \node[align=center] at (2.0,0.35)
    {solute first, then solvent\\\textbf{up to} the mark};
  \node[align=center,font=\sffamily\footnotesize] at (2.0,-0.55)
    {final \emph{solution} volume\\is exactly \SI{1.00}{\litre}};

  % WRONG WAY
  \begin{scope}[xshift=7.4cm]
    \node[font=\sffamily\small\bfseries] at (2.0,4.5) {what it is
      \textbf{not}};
    % fill BEFORE stroke -- the other way round the liquid paints over
    % the beaker walls and only the rim survives.
    \fill[apcgray!30] (0.55,1.35) rectangle (1.75,2.95);
    \draw[very thick] (0.55,3.2) rectangle (1.75,1.35);
    \node[font=\tiny] at (1.15,3.55) {\SI{1.00}{\litre} water};
    \node[font=\large] at (2.25,2.2) {$+$};
    \fill[apcgray!60] (3.0,1.9) circle (0.34);
    \node[font=\tiny] at (3.0,1.4) {solute};
    \draw[very thick,apcgray] (0.3,0.75) -- (3.6,3.6);
    \node[align=center] at (2.0,0.35)
      {a mole of solute added \emph{to}\\\SI{1.00}{\litre} of water};
    \node[align=center,font=\sffamily\footnotesize] at (2.0,-0.55)
      {the solute takes up room ---\\the volume is now
       \textbf{more} than \SI{1}{\litre}};
  \end{scope}
\end{tikzpicture}
\end{center}

\begin{workedexample}[I do: molarity from a mass]
\textbf{What is the molarity of a solution made by dissolving
\SI{25.0}{\gram} of \ce{NaCl} in enough water to make
\SI{500.}{\milli\litre} of solution?}

\textbf{Step 1 --- grams to moles.} You can never put grams into a
molarity; the definition demands moles.
\[ M(\ce{NaCl}) = 22.99 + 35.45 = \SI{58.44}{\gram\per\mole} \]
\[ n = \frac{\SI{25.0}{\gram}}{\SI{58.44}{\gram\per\mole}}
     = \SI{0.428}{\mole} \]

\textbf{Step 2 --- millilitres to litres.} \SI{500.}{\milli\litre}
$= \SI{0.500}{\litre}$.

\textbf{Step 3 --- divide.}
\[ M = \frac{\SI{0.428}{\mole}}{\SI{0.500}{\litre}}
     = \mathbf{\SI{0.856}{\Molar}} \]

\textbf{Two habits worth building now.} Convert to moles and litres
\emph{before} you divide, every time --- most lost marks here are unit
slips, not concept failures. And read the wording: ``dissolved in enough
water to make \SI{500.}{\milli\litre}'' is a solution volume;
``dissolved in \SI{500.}{\milli\litre} of water'' is a solvent volume,
and the two are not the same quantity.
\end{workedexample}

\begin{yourturn}[YOUR TURN 1 --- four questions]
\begin{enumerate}[leftmargin=*,itemsep=0.5em]
  \item Find the molarity of \SI{0.500}{\mole} of \ce{KCl} in
        \SI{2.00}{\litre} of solution. \workspace[1.6cm]
  \item Find the molarity of \SI{40.0}{\gram} of \ce{NaOH}
        ($M = \SI{40.00}{\gram\per\mole}$) in \SI{250.}{\milli\litre}
        of solution. \workspace[1.8cm]
  \item How many moles of solute are in \SI{50.0}{\milli\litre} of
        \SI{2.00}{\Molar} \ce{HCl}? \workspace[1.6cm]
  \item Why is molarity defined per litre of solution rather than per
        litre of solvent? \workspace[1.8cm]
\end{enumerate}
\selfcheck{(a) \SI{0.250}{\Molar} \quad (b) \SI{4.00}{\Molar} \quad
(c) \SI{0.100}{\mole} \quad (d) the solute itself occupies volume}
\keyonly{\begin{apcanswer}(a) $0.500 / 2.00 = \mathbf{\SI{0.250}{\Molar}}$.
(b) $n = 40.0/40.00 = \SI{1.00}{\mole}$; $1.00/0.250 =
\mathbf{\SI{4.00}{\Molar}}$. (c) $n = MV = 2.00 \times 0.0500 =
\mathbf{\SI{0.100}{\mole}}$. (d) Because the \textbf{solute takes up
space of its own}. Adding a mole of solute to exactly one litre of
solvent gives a total volume larger than one litre --- and by an amount
that differs from solute to solute, so it could not be predicted.
Defining molarity against the final solution volume makes it something
you can measure directly with a volumetric flask.\end{apcanswer}}
\end{yourturn}

% =====================================================================
\section*{\sffamily Ladder 2 \textbullet\ Dilution}
\zsec{11.1}

Adding solvent changes the volume and the concentration but
\textbf{not} the number of moles of solute. That single conserved
quantity is the whole calculation.

\[ \underbrace{M_{1}V_{1}}_{\text{moles before}}
   \;=\;
   \underbrace{M_{2}V_{2}}_{\text{moles after}} \]

\begin{center}
\begin{tikzpicture}[font=\sffamily\footnotesize]
  % concentrated
  \node[font=\sffamily\small\bfseries] at (1.3,3.5) {concentrated};
  \fill[apcgray!22] (0.5,0.4) rectangle (2.1,1.5);
  \draw[very thick] (0.5,2.6) -- (0.5,0.4) -- (2.1,0.4) -- (2.1,2.6);
  \foreach \x/\y in {0.75/0.65, 1.15/0.62, 1.55/0.7, 1.9/0.6,
                     0.7/1.0, 1.1/0.95, 1.5/1.05, 1.85/0.98,
                     0.9/1.32, 1.35/1.28, 1.75/1.35, 0.62/1.28} {
    \fill (\x,\y) circle (0.075);}
  \node at (1.3,-0.15) {12 particles};
  \node at (1.3,-0.6) {small volume};

  \draw[-{Stealth[length=3mm]},very thick] (2.6,1.3) -- (4.0,1.3);
  \node[above,align=center,font=\sffamily\footnotesize] at (3.3,1.45)
    {add\\solvent};

  % dilute
  \begin{scope}[xshift=4.5cm]
    \node[font=\sffamily\small\bfseries] at (1.3,3.5) {dilute};
    \fill[apcgray!22] (0.5,0.4) rectangle (2.1,2.35);
    \draw[very thick] (0.5,2.6) -- (0.5,0.4) -- (2.1,0.4) -- (2.1,2.6);
    \foreach \x/\y in {0.75/0.7, 1.45/0.62, 1.95/0.85,
                       0.62/1.25, 1.15/1.1, 1.8/1.35,
                       0.85/1.65, 1.55/1.75, 1.98/1.6,
                       0.68/2.05, 1.25/2.15, 1.85/2.2} {
      \fill (\x,\y) circle (0.075);}
    \node at (1.3,-0.15) {\textbf{still} 12 particles};
    \node at (1.3,-0.6) {larger volume};
  \end{scope}

  \node[anchor=west,align=left,font=\sffamily\footnotesize]
    at (7.5,1.55)
    {the \textbf{number} of solute particles\\
     never changes --- only how much\\
     room they are spread through};
\end{tikzpicture}
\end{center}

\begin{workedexample}[I do: two directions of the same equation]
\textbf{(a) What is the concentration when \SI{50.0}{\milli\litre} of
\SI{6.00}{\Molar} \ce{HCl} is diluted to \SI{250.}{\milli\litre}?}

Rearrange for what you want, then substitute:
\[ M_{2} = \frac{M_{1}V_{1}}{V_{2}}
         = \frac{6.00 \times 50.0}{250.}
         = \mathbf{\SI{1.20}{\Molar}} \]

\textbf{Note the units did not need converting.} $V_{1}$ and $V_{2}$
appear as a ratio, so as long as both are in the same unit the
millilitres cancel. This is the one place in solution work where you may
safely leave millilitres alone.

\textbf{(b) What volume of \SI{6.00}{\Molar} \ce{HCl} is needed to make
\SI{2.00}{\litre} of \SI{0.150}{\Molar} \ce{HCl}?}
\[ V_{1} = \frac{M_{2}V_{2}}{M_{1}}
         = \frac{0.150 \times 2.00}{6.00}
         = \SI{0.0500}{\litre} = \mathbf{\SI{50.0}{\milli\litre}} \]

\textbf{And the safety rule that goes with it.} You would measure that
\SI{50.0}{\milli\litre} of acid and add it \emph{to} water already in
the flask --- never water to concentrated acid. Dissolving concentrated
acid is strongly exothermic, and adding water to it can boil the surface
and spit acid out of the container.
\end{workedexample}

\begin{yourturn}[YOUR TURN 2 --- four questions]
\begin{enumerate}[leftmargin=*,itemsep=0.5em]
  \item \SI{25.0}{\milli\litre} of \SI{12.0}{\Molar} \ce{HCl} is diluted
        to \SI{500.}{\milli\litre}. Find the new molarity.
        \workspace[1.6cm]
  \item What volume of \SI{2.00}{\Molar} stock makes
        \SI{100.}{\milli\litre} of \SI{0.400}{\Molar} solution?
        \workspace[1.6cm]
  \item \SI{100.}{\milli\litre} of \SI{0.50}{\Molar} \ce{NaCl} has
        \SI{100.}{\milli\litre} of water added. Find the new molarity.
        \workspace[1.8cm]
  \item Does diluting a solution change the number of moles of solute?
        Explain. \workspace[1.6cm]
\end{enumerate}
\selfcheck{(a) \SI{0.600}{\Molar} \quad (b) \SI{20.0}{\milli\litre}
\quad (c) \SI{0.25}{\Molar} \quad (d) no --- only the volume changes}
\keyonly{\begin{apcanswer}(a) $M_{2} = (12.0 \times 25.0)/500. =
\mathbf{\SI{0.600}{\Molar}}$. (b) $V_{1} = (0.400 \times 100.)/2.00 =
\mathbf{\SI{20.0}{\milli\litre}}$. (c) The final volume is
$100. + 100. = \SI{200.}{\milli\litre}$, so $M_{2} = (0.50 \times
100.)/200. = \mathbf{\SI{0.25}{\Molar}}$ --- the concentration halves
because the volume doubled. (d) \textbf{No.} Dilution adds only solvent,
so every solute particle that was there before is still there
afterwards. What changes is the volume they are distributed
through, and therefore the concentration. That conservation is exactly
what $M_{1}V_{1} = M_{2}V_{2}$ states.\end{apcanswer}}
\end{yourturn}

% =====================================================================
\section*{\sffamily Ladder 3 \textbullet\ Mass percent, ppm and ppb}
\zsec{11.1}

When the solute is very dilute, molarity produces awkward numbers.
These units exist to make small amounts readable.

\begin{center}
\begin{tikzpicture}[font=\sffamily\footnotesize]
  \draw[-{Stealth[length=3mm]},very thick] (0.3,2.2) -- (11.6,2.2);
  \foreach \x/\u/\f/\e in {%
      1.3/{mass \%}/{parts per hundred}/{salt in seawater, 3.5\%},
      5.0/{ppm}/{parts per million}/{fluoride in tap water, 1 ppm},
      8.9/{ppb}/{parts per billion}/{lead limit, 15 ppb}} {
    \fill (\x,2.2) circle (0.09);
    \node[font=\sffamily\bfseries] at (\x,2.75) {\u};
    \node at (\x,1.72) {\f};
    \node[font=\sffamily\footnotesize\itshape] at (\x,1.25) {\e};}
  % Raised to 3.35: at 2.65 this ran straight into the "ppb" label,
  % printing as "ppb_more dilute".
  \node[anchor=east,font=\sffamily\footnotesize] at (11.6,3.35)
    {more dilute $\longrightarrow$};
  \draw[apcgray] (0.3,0.72) -- (11.6,0.72);
  \node[anchor=west,font=\sffamily\footnotesize] at (0.4,0.3)
    {each is a \textbf{mass ratio} $\times$ a scaling factor:
     $\times 10^{2}$, $\times 10^{6}$, $\times 10^{9}$};
\end{tikzpicture}
\end{center}

\[ \text{mass \%} =
   \frac{\text{mass of solute}}{\text{mass of \textbf{solution}}}
   \times 100 \]

\begin{apctrap}
The denominator is the mass of the \textbf{solution}, meaning
solute $+$ solvent --- not the mass of the solvent alone. Dissolving
\SI{15}{\gram} in \SI{135}{\gram} of water gives a
\SI{150}{\gram} solution, so the answer is $15/150$, not $15/135$. The
same total-versus-part trap as molarity, one ladder later.
\end{apctrap}

\begin{workedexample}[I do: mass percent and ppm]
\textbf{(a) \SI{15.0}{\gram} of \ce{NaCl} is dissolved in
\SI{135.0}{\gram} of water. Find the mass percent.}

Total solution mass $= 15.0 + 135.0 = \SI{150.0}{\gram}$.
\[ \text{mass \%} = \frac{15.0}{150.0} \times 100
                  = \mathbf{10.0\%} \]

\textbf{(b) A water sample contains \SI{2.5}{\milli\gram} of fluoride
per \SI{1.0}{\kilo\gram} of water. Express this in ppm.}

Parts per million is the same ratio scaled by $10^{6}$. Put both masses
in the same unit first: \SI{2.5}{\milli\gram} $=
\SI{2.5e-3}{\gram}$ and \SI{1.0}{\kilo\gram} $= \SI{1.0e3}{\gram}$.
\[ \frac{\num{2.5e-3}}{\num{1.0e3}} \times 10^{6}
   = \mathbf{2.5 \text{ ppm}} \]

\textbf{The shortcut worth memorising.} For dilute \emph{aqueous}
solutions the density is essentially \SI{1.00}{\gram\per\milli\litre},
so \SI{1}{\kilo\gram} of solution is \SI{1}{\litre} and
\[ 1 \text{ ppm} \;\approx\; \SI{1}{\milli\gram\per\litre} \]
That is why water-quality reports quote mg/L and ppm interchangeably.
It is an approximation that holds only because the solution is mostly
water --- do not use it for concentrated solutions or non-aqueous
solvents.
\end{workedexample}

\begin{yourturn}[YOUR TURN 3 --- four questions]
\begin{enumerate}[leftmargin=*,itemsep=0.5em]
  \item Find the mass percent of \SI{5.0}{\gram} of sugar in
        \SI{95.0}{\gram} of water. \workspace[1.6cm]
  \item A solution is 2.0\% by mass. How much solute is in
        \SI{250}{\gram} of it? \workspace[1.6cm]
  \item Convert \SI{0.0025}{\gram} of solute per \SI{1000}{\gram} of
        solution to ppm. \workspace[1.6cm]
  \item Why is the denominator the mass of solution rather than of
        solvent? \workspace[1.6cm]
\end{enumerate}
\selfcheck{(a) 5.0\% \quad (b) \SI{5.0}{\gram} \quad (c) 2.5 ppm \quad
(d) it is a fraction of the whole mixture}
\keyonly{\begin{apcanswer}(a) Total $= 5.0 + 95.0 = \SI{100.0}{\gram}$,
so $5.0/100.0 \times 100 = \mathbf{5.0\%}$. (b) $0.020 \times 250 =
\mathbf{\SI{5.0}{\gram}}$. (c) $(0.0025/1000) \times 10^{6} =
\mathbf{2.5}$ \textbf{ppm}. (d) Because these units express what
\textbf{fraction of the whole mixture} is solute. A 10\% solution means
ten grams in every hundred grams of \emph{solution}; if the denominator
were the solvent, two solutions with the same label would contain
different amounts of solute per gram of liquid you actually pour, and
the number would not mean anything consistent.\end{apcanswer}}
\end{yourturn}

% =====================================================================
\section*{\sffamily Ladder 4 \textbullet\ Mole fraction and molality}
\zsec{11.1}

Four concentration units, and what separates them is entirely
\textbf{what sits in the denominator}. Get that right and all four are
one idea.

\begin{center}
\begin{tikzpicture}[font=\sffamily\footnotesize]
  \node[font=\sffamily\small\bfseries] at (5.6,3.5)
    {same numerator idea --- four different denominators};
  \foreach \y/\n/\sym/\num/\den in {%
      2.7/{molarity}/{$M$}/{mol solute}/{\textbf{litres of solution}},
      2.0/{molality}/{$m$}/{mol solute}/{\textbf{kg of solvent}},
      1.3/{mole fraction}/{$\chi$}/{mol solute}/{\textbf{total mol,
        solute $+$ solvent}},
      0.6/{mass percent}/{---}/{mass solute}/{\textbf{mass of
        solution}}} {
    \node[anchor=west,font=\sffamily\bfseries] at (0.2,\y) {\n};
    \node[anchor=west] at (2.6,\y) {\sym};
    \node[anchor=west] at (3.3,\y) {\num};
    \node[anchor=west] at (5.6,\y) {$\div$};
    \node[anchor=west] at (6.2,\y) {\den};}
  \draw[apcgray] (0.1,3.1) -- (11.4,3.1);
  \draw[apcgray] (0.1,0.25) -- (11.4,0.25);
\end{tikzpicture}
\end{center}

\begin{apcnote}[Where molality actually gets used --- read this once]
Molality exists mainly to support \textbf{colligative properties}, and
those are \S11.4--11.7, which this course skips. So learn the definition
and learn how it differs from molarity --- that distinction is fair game
--- but do not expect to apply it. Two reasons it is defined against
solvent \emph{mass} rather than solution \emph{volume}: mass does not
change with temperature, whereas volume does, so molality is
temperature-independent where molarity is not.
\end{apcnote}

\begin{workedexample}[I do: mole fraction, then molality]
\textbf{(a) A solution contains \SI{46.0}{\gram} of ethanol
(\ce{C2H5OH}, $M = \SI{46.07}{\gram\per\mole}$) and \SI{54.0}{\gram}
of water ($M = \SI{18.02}{\gram\per\mole}$). Find the mole fraction of
ethanol.}

\textbf{Everything must be in moles} --- a mole fraction cannot be built
from grams.
\[ n_{\text{ethanol}} = \frac{46.0}{46.07} = \SI{0.998}{\mole}
   \qquad
   n_{\text{water}} = \frac{54.0}{18.02} = \SI{3.00}{\mole} \]
\[ \chi_{\text{ethanol}}
   = \frac{0.998}{0.998 + 3.00}
   = \frac{0.998}{4.00}
   = \mathbf{0.250} \]

\textbf{Check it the free way.} Mole fractions of everything present
must sum to exactly 1, so $\chi_{\text{water}} = 1 - 0.250 = 0.750$.
If your fractions do not sum to 1, you have made an arithmetic error and
the check has just caught it.

\textbf{Note what happened to the masses.} Nearly equal masses
(\SI{46.0}{} and \SI{54.0}{\gram}) gave wildly unequal mole counts
(\SI{0.998}{} and \SI{3.00}{\mole}), because a water molecule is so much
lighter. Mass ratios and mole ratios are not interchangeable.

\textbf{(b) Find the molality of \SI{18.0}{\gram} of glucose
(\ce{C6H12O6}, $M = \SI{180.16}{\gram\per\mole}$) dissolved in
\SI{250.}{\gram} of water.}
\[ n = \frac{18.0}{180.16} = \SI{0.0999}{\mole}
   \qquad
   m = \frac{0.0999}{\SI{0.250}{\kilo\gram}}
     = \mathbf{\SI{0.400}{\molal}} \]
The denominator is \SI{0.250}{\kilo\gram} of \emph{water}, not of
solution --- for molality the solute mass never enters the bottom.
\end{workedexample}

\begin{yourturn}[YOUR TURN 4 --- four questions]
\begin{enumerate}[leftmargin=*,itemsep=0.5em]
  \item A mixture holds \SI{2.00}{\mole} of \ce{A} and
        \SI{6.00}{\mole} of \ce{B}. Find $\chi_{\ce{A}}$.
        \workspace[1.6cm]
  \item If $\chi_{\ce{A}} = 0.35$ in a two-component mixture, what is
        $\chi_{\ce{B}}$? \workspace[1.4cm]
  \item Find the molality of \SI{0.500}{\mole} of solute in
        \SI{2.00}{\kilo\gram} of solvent. \workspace[1.6cm]
  \item State the one difference between molarity and molality, and say
        why molality does not change with temperature.
        \workspace[1.8cm]
\end{enumerate}
\selfcheck{(a) 0.250 \quad (b) 0.65 \quad (c) \SI{0.250}{\molal} \quad
(d) litres of solution vs kg of solvent; mass is
temperature-independent}
\keyonly{\begin{apcanswer}(a) $2.00/(2.00+6.00) = \mathbf{0.250}$.
(b) Mole fractions sum to 1, so $1 - 0.35 = \mathbf{0.65}$.
(c) $0.500/2.00 = \mathbf{\SI{0.250}{\molal}}$. (d) Molarity divides by
\textbf{litres of solution}; molality divides by \textbf{kilograms of
solvent}. Molality is temperature-independent because \textbf{mass does
not change with temperature} --- but volume does, since liquids expand
when heated. Warm a solution and its molarity falls slightly (same moles,
larger volume) while its molality is unchanged.\end{apcanswer}}
\end{yourturn}

% =====================================================================
\section*{\sffamily Ladder 5 \textbullet\ The energetics of dissolving}
\zsec{11.2}

Dissolving is not one event but three, and only the third releases
energy. Whether a solution forms --- and whether the beaker gets warm or
cold --- is the sum.

\begin{center}
\begin{tikzpicture}[font=\sffamily\footnotesize]
  \foreach \i/\lab/\desc/\sign in {%
      0/{Step 1}/{separate the\\\textbf{solute} particles}/{endo, $+$},
      4.0/{Step 2}/{separate the\\\textbf{solvent} particles}/{endo, $+$},
      8.0/{Step 3}/{let solute and solvent\\\textbf{attract} each
        other}/{exo, $-$}} {
    \begin{scope}[xshift=\i cm]
      \node[font=\sffamily\small\bfseries] at (1.6,3.35) {\lab};
      \node[align=center] at (1.6,2.7) {\desc};
      % Sign sits BELOW the art at 1.15.  At 1.55 it was inside the
      % particle sketches -- the step-1 arrow ran straight through
      % "endo, +" and step 3's solvation shell through "exo, -".
      \node[font=\sffamily\bfseries] at (1.6,1.15) {\sign};
    \end{scope}}
  % step 1 art: clump -> spread
  \foreach \p in {(0.75,2.0),(1.05,2.0),(0.75,1.75),(1.05,1.75)} {
    \fill[apcgray!65] \p circle (0.1);}
  \draw[-{Stealth[length=2mm]}] (1.35,1.88) -- (1.75,1.88);
  \foreach \p in {(2.0,2.1),(2.4,1.95),(2.15,1.65),(2.55,2.2)} {
    \fill[apcgray!65] \p circle (0.1);}
  % step 2 art
  \begin{scope}[xshift=4.0cm]
    \foreach \p in {(0.75,2.0),(1.05,2.0),(0.75,1.75),(1.05,1.75)} {
      \fill[apcgray!25] \p circle (0.1);
      \draw[apcgray] \p circle (0.1);}
    \draw[-{Stealth[length=2mm]}] (1.35,1.88) -- (1.75,1.88);
    \foreach \p in {(2.0,2.1),(2.4,1.95),(2.15,1.65),(2.55,2.2)} {
      \fill[apcgray!25] \p circle (0.1);
      \draw[apcgray] \p circle (0.1);}
  \end{scope}
  % step 3 art: solute surrounded by solvent
  \begin{scope}[xshift=8.0cm]
    \fill[apcgray!65] (1.6,1.95) circle (0.15);
    \foreach \a in {0,72,144,216,288} {
      \fill[apcgray!25] ([shift={(1.6,1.95)}]\a:0.36) circle (0.11);
      \draw[apcgray] ([shift={(1.6,1.95)}]\a:0.36) circle (0.11);}
  \end{scope}
  \draw[apcgray] (0.2,0.72) -- (11.2,0.72);
  \node[font=\sffamily\small] at (5.7,0.28)
    {$\Delta H_{\text{soln}} \;=\; \Delta H_{1} \;+\; \Delta H_{2}
     \;+\; \Delta H_{3}$};
\end{tikzpicture}
\end{center}

\begin{center}
\begin{tikzpicture}[font=\sffamily\footnotesize]
  % EXOTHERMIC
  \node[font=\sffamily\small\bfseries] at (2.1,3.5) {exothermic
    ($\Delta H_{\text{soln}} < 0$)};
  \draw[-{Stealth[length=2.5mm]},thick] (0.35,0.4) -- (0.35,3.0)
    node[above,font=\tiny] {$H$};
  \draw[very thick] (0.8,2.5) -- (2.1,2.5);
  \node[anchor=west,font=\tiny] at (2.2,2.5) {separated};
  \draw[very thick] (0.8,1.75) -- (2.1,1.75);
  \node[anchor=west,font=\tiny] at (2.2,1.75) {start};
  \draw[very thick] (0.8,0.85) -- (2.1,0.85);
  \node[anchor=west,font=\tiny] at (2.2,0.85) {solution};
  \draw[-{Stealth[length=2mm]},thick] (1.45,1.75) -- (1.45,0.85);
  \node[align=center] at (2.1,-0.4)
    {step 3 releases more than\\steps 1 and 2 cost ---
     \textbf{warms up}};

  % ENDOTHERMIC
  \begin{scope}[xshift=6.6cm]
    \node[font=\sffamily\small\bfseries] at (2.1,3.5) {endothermic
      ($\Delta H_{\text{soln}} > 0$)};
    \draw[-{Stealth[length=2.5mm]},thick] (0.35,0.4) -- (0.35,3.0)
      node[above,font=\tiny] {$H$};
    \draw[very thick] (0.8,2.5) -- (2.1,2.5);
    \node[anchor=west,font=\tiny] at (2.2,2.5) {separated};
    \draw[very thick] (0.8,0.85) -- (2.1,0.85);
    \node[anchor=west,font=\tiny] at (2.2,0.85) {start};
    \draw[very thick] (0.8,1.6) -- (2.1,1.6);
    \node[anchor=west,font=\tiny] at (2.2,1.6) {solution};
    \draw[-{Stealth[length=2mm]},thick] (1.45,0.85) -- (1.45,1.6);
    \node[align=center] at (2.1,-0.4)
      {step 3 releases less than\\steps 1 and 2 cost ---
       \textbf{cools down}};
  \end{scope}
\end{tikzpicture}
\end{center}

\begin{workedexample}[I do: why \ce{NaCl} dissolves at all]
\textbf{Breaking the \ce{NaCl} lattice costs
\SI{+786}{\kilo\joule\per\mole}. Hydrating the resulting ions releases
\SI{-783}{\kilo\joule\per\mole}. Find $\Delta H_{\text{soln}}$ and
comment.}

\[ \Delta H_{\text{soln}} = (+786) + (-783)
   = \mathbf{\SI{+3}{\kilo\joule\per\mole}} \]

\textbf{So dissolving salt is slightly endothermic} --- a beaker of
water cools very slightly as salt dissolves. And yet \ce{NaCl} is
famously soluble. Why does an uphill process happen at all?

\textbf{Because enthalpy is not the only driver.} A dissolved solute is
enormously more disordered than a crystal: the ions go from fixed
lattice sites to wandering freely through the liquid. That increase in
\textbf{entropy} is what drives dissolving here, and it is large enough
to overwhelm a \SI{3}{\kilo\joule\per\mole} enthalpy penalty.

\textbf{Look at the two numbers again.} \SI{786}{} and \SI{783}{} ---
each enormous, and their difference tiny. Solubility is decided by the
small residue left when two large opposing quantities nearly cancel,
which is why it is so hard to predict from first principles and why
chemists rely on measured solubility data instead.

\textbf{Two contrasting cases worth knowing.}
\ce{NH4NO3} has $\Delta H_{\text{soln}} =
\SI{+25.7}{\kilo\joule\per\mole}$ --- strongly endothermic, and the
active ingredient of an instant \emph{cold} pack.
\ce{CaCl2} has $\Delta H_{\text{soln}} =
\SI{-82.8}{\kilo\joule\per\mole}$ --- strongly exothermic, and used in
instant \emph{hot} packs and to melt road ice.
\end{workedexample}

\begin{yourturn}[YOUR TURN 5 --- four questions]
\begin{enumerate}[leftmargin=*,itemsep=0.5em]
  \item Which of the three steps is exothermic, and why?
        \workspace[1.6cm]
  \item Lattice breaking costs \SI{+700}{\kilo\joule\per\mole} and
        hydration releases \SI{-750}{\kilo\joule\per\mole}. Find
        $\Delta H_{\text{soln}}$ and say whether the beaker warms or
        cools. \workspace[1.8cm]
  \item A cold pack contains \ce{NH4NO3}. What is the sign of its
        $\Delta H_{\text{soln}}$? \workspace[1.4cm]
  \item \ce{NaCl} dissolving is endothermic, yet it dissolves readily.
        Explain. \workspace[1.8cm]
\end{enumerate}
\selfcheck{(a) step 3 --- attractions form \quad (b)
\SI{-50}{\kilo\joule\per\mole}, warms \quad (c) positive \quad (d)
entropy drives it}
\keyonly{\begin{apcanswer}(a) \textbf{Step 3}, in which solute and
solvent particles attract one another. Forming an attraction always
releases energy; breaking one always costs energy, which is why steps 1
and 2 are both endothermic. (b) $(+700) + (-750) =
\mathbf{\SI{-50}{\kilo\joule\per\mole}}$ --- negative, so the process is
exothermic and the beaker \textbf{warms}. (c) \textbf{Positive}
(endothermic). The pack gets cold because dissolving absorbs energy from
its surroundings, which include your skin. (d) Because dissolving is
driven by \textbf{entropy} as well as enthalpy. The ions move from fixed
positions in an ordered lattice to moving freely through the solution, an
enormous increase in disorder, and that is more than enough to overcome
an enthalpy cost of only about
\SI{3}{\kilo\joule\per\mole}.\end{apcanswer}}
\end{yourturn}

% =====================================================================
\section*{\sffamily Ladder 6 \textbullet\ Like dissolves like}
\zsec{11.2--11.3}

Step 3 of Ladder 5 only pays for steps 1 and 2 when the new
solute--solvent attractions resemble the ones that were broken. That is
the whole content of the slogan.

\begin{center}
\begin{tikzpicture}[font=\sffamily\footnotesize]
  \node[font=\sffamily\small\bfseries] at (4.4,4.3) {will it dissolve?};
  % column headings
  \node[font=\sffamily\bfseries,align=center] at (3.3,3.5)
    {\textbf{polar} solvent\\(water)};
  \node[font=\sffamily\bfseries,align=center] at (6.6,3.5)
    {\textbf{nonpolar} solvent\\(hexane)};
  % row headings
  \node[anchor=east,font=\sffamily\bfseries,align=right] at (1.9,2.35)
    {\textbf{polar} or ionic solute\\(\ce{NaCl}, sugar)};
  \node[anchor=east,font=\sffamily\bfseries,align=right] at (1.9,1.0)
    {\textbf{nonpolar} solute\\(\ce{I2}, oil)};
  % grid
  \draw[apcgray] (2.1,3.0) -- (8.2,3.0);
  \draw[apcgray] (2.1,1.68) -- (8.2,1.68);
  \draw[apcgray] (4.95,3.15) -- (4.95,0.35);
  % cells
  \node[font=\Large] at (3.3,2.35) {\checkmark};
  \node[font=\Large] at (6.6,2.35) {$\times$};
  \node[font=\Large] at (3.3,1.0) {$\times$};
  \node[font=\Large] at (6.6,1.0) {\checkmark};
  \node[font=\sffamily\footnotesize] at (5.15,-0.3)
    {the diagonal fails --- \emph{unlike} does \textbf{not} dissolve
     \emph{unlike}};
\end{tikzpicture}
\end{center}

\begin{workedexample}[I do: why oil and water refuse]
\textbf{Water is strongly hydrogen bonded. Oil is a nonpolar
hydrocarbon with dispersion forces only. Work through the three steps.}

\textbf{Step 1 --- separate the oil.} Cheap. Only weak dispersion forces
have to be overcome, so this costs little.

\textbf{Step 2 --- separate the water.} \emph{Expensive.} Every water
molecule is hydrogen bonded to its neighbours, and opening a cavity big
enough for an oil molecule means breaking a great many of those bonds.

\textbf{Step 3 --- oil attracts water.} Cheap and useless. An oil
molecule cannot hydrogen-bond, so all it offers water in return is weak
dispersion --- nowhere near enough to repay what step 2 cost.

\textbf{The books do not balance,} so the mixture separates and the
water re-forms its hydrogen-bond network with itself.

\textbf{Now reverse the pairing and watch it work.} Ethanol,
\ce{CH3CH2OH}, has an \ce{-OH} group and hydrogen-bonds to water just
as water does to itself. Step 2 is still expensive, but step 3 now
repays it in the same currency --- so ethanol and water are
\textbf{miscible} in all proportions.

\textbf{Say it precisely on an exam.} Not ``like dissolves like'', which
earns nothing, but: ``the solute--solvent attractions are comparable in
strength to the solvent--solvent attractions that must be broken.''
Naming the forces is what scores.
\end{workedexample}

\begin{yourturn}[YOUR TURN 6 --- four questions]
\begin{enumerate}[leftmargin=*,itemsep=0.5em]
  \item Will \ce{I2} dissolve better in water or in \ce{CCl4}? Why?
        \workspace[1.8cm]
  \item Will \ce{NaCl} dissolve in hexane? Why not?
        \workspace[1.8cm]
  \item Why is ethanol miscible with water in all proportions?
        \workspace[1.8cm]
  \item Which step of the cycle is the expensive one when a nonpolar
        solute meets water? \workspace[1.6cm]
\end{enumerate}
\selfcheck{(a) \ce{CCl4} --- both nonpolar \quad (b) no --- hexane
cannot solvate ions \quad (c) it hydrogen-bonds \quad (d) step 2}
\keyonly{\begin{apcanswer}(a) \textbf{\ce{CCl4}}. Both \ce{I2} and
\ce{CCl4} are nonpolar with dispersion forces only, so the attractions
broken and formed are of the same kind and comparable strength. In
water, \ce{I2} could not repay the cost of breaking water's hydrogen
bonds. (b) \textbf{No.} Separating the \ce{NaCl} lattice costs an
enormous amount of energy, and hexane is nonpolar --- it can offer the
ions only weak dispersion forces, nothing like the ion--dipole
attractions water provides. Step 3 cannot pay for step 1.
(c) Because ethanol's \ce{-OH} group lets it \textbf{hydrogen-bond to
water}, so the ethanol--water attractions formed are as strong as the
water--water attractions broken. (d) \textbf{Step 2}, separating the
solvent. Water's hydrogen-bond network is strong, and opening a space
for the solute means breaking a lot of it --- a cost a nonpolar solute
has no way to repay.\end{apcanswer}}
\end{yourturn}

% =====================================================================
\section*{\sffamily Ladder 7 \textbullet\ Structure: chains, heads and
tails}
\zsec{11.3}

Most real molecules are neither purely polar nor purely nonpolar. They
have a \textbf{hydrophilic} part that water likes and a
\textbf{hydrophobic} part it does not, and solubility follows whichever
wins.

\begin{center}
\begin{tikzpicture}[font=\sffamily\footnotesize]
  \node[font=\sffamily\small\bfseries] at (5.6,3.6)
    {straight-chain alcohols: one \ce{-OH}, longer and longer tails};
  \node[font=\sffamily\footnotesize] at (5.6,3.15)
    {dark circle $=$ the \ce{-OH} head \textbullet\ pale circles $=$
     carbons of the tail};
  \foreach \i/\nm/\n/\sol in {%
      0/{methanol}/1/{miscible},
      2.4/{ethanol}/2/{miscible},
      4.8/{1-butanol}/4/{7.9},
      7.2/{1-hexanol}/6/{0.6},
      9.6/{1-octanol}/8/{0.054}} {
    \begin{scope}[xshift=\i cm]
      \node[font=\sffamily\footnotesize] at (0.9,2.85) {\nm};
      % OH head.  No lettering inside: "OH" even at \tiny is wider than
      % the circle and rendered as an illegible smudge.  The key below
      % the title identifies it instead.
      \fill[apcgray!75] (0.25,2.2) circle (0.15);
      % tail
      \foreach \k in {1,...,\n} {
        \fill[apcgray!30] (0.25+\k*0.16,2.2) circle (0.085);}
      \node[font=\sffamily\footnotesize] at (0.9,1.45) {\sol};
    \end{scope}}
  \node[anchor=east,font=\sffamily\footnotesize] at (-0.1,2.2)
    {structure};
  \node[anchor=east,font=\sffamily\footnotesize,align=right]
    at (-0.1,1.45) {solubility\\g/100 g water};
  \draw[-{Stealth[length=3mm]},thick,apcgray] (0.3,0.85) -- (10.9,0.85);
  \node[font=\sffamily\footnotesize] at (5.6,0.42)
    {the hydrophobic tail grows, the single hydrophilic head does not
     --- solubility collapses};
\end{tikzpicture}
\end{center}

\begin{center}
\begin{tikzpicture}[font=\sffamily\footnotesize]
  \node[font=\sffamily\small\bfseries] at (2.4,3.7) {a soap molecule};
  \fill[apcgray!70] (0.6,2.4) circle (0.24);
  \draw[very thick,decorate,decoration={snake,amplitude=1.2pt,
    segment length=5pt}] (0.86,2.4) -- (3.9,2.4);
  \node[align=center,font=\sffamily\footnotesize] at (0.6,1.75)
    {ionic \textbf{head}\\hydrophilic};
  \node[align=center,font=\sffamily\footnotesize] at (2.6,1.75)
    {hydrocarbon \textbf{tail}\\hydrophobic};

  % micelle
  \begin{scope}[xshift=7.6cm]
    \node[font=\sffamily\small\bfseries] at (1.9,3.7) {a micelle};
    \fill[apcgray!12] (1.9,2.15) circle (1.35);
    \foreach \a in {0,30,...,330} {
      \fill[apcgray!70] ([shift={(1.9,2.15)}]\a:1.28) circle (0.13);
      \draw[thick] ([shift={(1.9,2.15)}]\a:1.13)
                -- ([shift={(1.9,2.15)}]\a:0.42);}
    \node[font=\tiny,align=center] at (1.9,2.15) {grease\\inside};
    \node[align=center,font=\sffamily\footnotesize] at (1.9,0.35)
      {heads face the water, tails\\trap the grease --- how soap works};
  \end{scope}
\end{tikzpicture}
\end{center}

\begin{workedexample}[I do: reading a solubility collapse]
\textbf{Methanol is miscible with water in all proportions.
1-octanol dissolves to only \SI{0.054}{\gram} per \SI{100}{\gram} of
water --- a difference of orders of magnitude. Both have exactly one
\ce{-OH} group.}

\textbf{What stayed the same.} Each molecule has a single \ce{-OH}, so
each can form roughly the same number of hydrogen bonds to water. The
hydrophilic contribution is essentially constant across the series.

\textbf{What changed.} The hydrocarbon tail grew from one carbon to
eight. That tail cannot hydrogen-bond, so it forces water to break its
own hydrogen bonds and open a cavity, paying step 2's cost with nothing
in return. The longer the tail, the bigger the unpaid bill.

\textbf{So the ratio decides.} In methanol the \ce{-OH} dominates the
molecule and it behaves like water. By octanol the tail dominates and
the molecule behaves like a hydrocarbon. The crossover sits around four
carbons, which is why 1-butanol is only partly soluble
(\SI{7.9}{\gram}/\SI{100}{\gram}) --- it is the molecule where the two
halves are most nearly matched.

\textbf{Where this shows up outside chemistry class.} Vitamin C is
studded with \ce{-OH} groups and is water soluble, so the body cannot
store it and you need it regularly. Vitamin A is mostly hydrocarbon and
is fat soluble, so it accumulates in body fat --- which is why an
overdose of vitamin A is possible and an overdose of vitamin C is
largely just expensive.
\end{workedexample}

\begin{yourturn}[YOUR TURN 7 --- four questions]
\begin{enumerate}[leftmargin=*,itemsep=0.5em]
  \item Which is more soluble in water, 1-propanol or 1-heptanol? Why?
        \workspace[1.8cm]
  \item Which part of a soap molecule dissolves grease?
        \workspace[1.6cm]
  \item Why is vitamin A stored in body fat but vitamin C is not?
        \workspace[1.8cm]
  \item Would you expect ethanoic acid (\ce{CH3COOH}) to be water
        soluble? Explain. \workspace[1.8cm]
\end{enumerate}
\selfcheck{(a) 1-propanol --- shorter tail \quad (b) the hydrocarbon
tail \quad (c) A is hydrophobic, C is hydrophilic \quad (d) yes ---
short chain, polar \ce{-COOH}}
\keyonly{\begin{apcanswer}(a) \textbf{1-propanol.} Both have one
\ce{-OH}, so both hydrogen-bond to water about equally, but propanol's
three-carbon tail is far shorter than heptanol's seven. Less hydrophobic
tail means less of water's hydrogen-bond network has to be broken for
nothing in return. (b) The \textbf{hydrocarbon tail}. It is nonpolar, so
it mixes with nonpolar grease by dispersion forces, while the ionic head
stays in the water --- letting one molecule bridge two phases that
otherwise refuse to mix. (c) \textbf{Vitamin A} is mostly hydrocarbon
and therefore hydrophobic, so it dissolves in fatty tissue and
accumulates there. \textbf{Vitamin C} carries several \ce{-OH} groups,
is hydrophilic, dissolves in the body's water and is excreted --- so it
must be replaced regularly. (d) \textbf{Yes.} It has a very short
(one-carbon) hydrophobic tail and a \ce{-COOH} group that both donates
and accepts hydrogen bonds. The hydrophilic part dominates, and ethanoic
acid is in fact miscible with water --- as anyone who has used vinegar
has observed.\end{apcanswer}}
\end{yourturn}

% =====================================================================
\section*{\sffamily Ladder 8 \textbullet\ Saturated, unsaturated,
supersaturated}
\zsec{11.3}

Solubility is a \emph{limit}, and a solution is classified by where it
sits relative to that limit.

\begin{center}
\begin{tikzpicture}[font=\sffamily\footnotesize]
  \foreach \i/\nm/\desc in {%
      0/{unsaturated}/{below the limit\\more \emph{would} dissolve},
      4.2/{saturated}/{at the limit\\undissolved solid remains},
      8.4/{supersaturated}/{\emph{above} the limit\\unstable}} {
    \begin{scope}[xshift=\i cm]
      \node[font=\sffamily\small\bfseries] at (1.5,3.5) {\nm};
      \fill[apcgray!22] (0.55,0.5) rectangle (2.45,2.35);
      \draw[very thick] (0.55,2.9) -- (0.55,0.5) -- (2.45,0.5)
        -- (2.45,2.9);
      \node[align=center] at (1.5,-0.35) {\desc};
    \end{scope}}
  % unsaturated: few dissolved particles, no solid
  \foreach \p in {(0.85,0.9),(1.4,1.3),(1.95,0.85),(1.1,1.9),
                  (2.15,1.7),(0.75,1.5)} {
    \fill \p circle (0.075);}
  % saturated: many dissolved + solid at the bottom
  \begin{scope}[xshift=4.2cm]
    \foreach \p in {(0.8,1.0),(1.25,1.45),(1.7,0.95),(2.1,1.35),
                    (0.95,1.95),(1.55,2.1),(2.2,1.85),(0.7,1.6)} {
      \fill \p circle (0.075);}
    \fill[apcgray!75] (0.9,0.5) rectangle (1.35,0.68);
    \fill[apcgray!75] (1.5,0.5) rectangle (2.0,0.64);
  \end{scope}
  % supersaturated: very many dissolved, no solid
  \begin{scope}[xshift=8.4cm]
    \foreach \p in {(0.75,0.85),(1.15,1.05),(1.6,0.85),(2.05,1.05),
                    (0.9,1.4),(1.35,1.55),(1.8,1.35),(2.2,1.6),
                    (0.75,1.9),(1.25,2.1),(1.75,1.95),(2.15,2.2)} {
      \fill \p circle (0.075);}
  \end{scope}
\end{tikzpicture}
\end{center}

\begin{apctrap}
A \textbf{supersaturated} solution holds more solute than it should be
able to. It is made by saturating a solution hot and then cooling it
slowly and undisturbed, so the excess has no site on which to
crystallise. Drop in a single seed crystal --- or scratch the glass ---
and the whole excess crystallises at once, releasing its heat of
crystallisation. That is exactly how a reusable sodium acetate hand
warmer works, and why clicking its metal disc turns it solid and hot in
seconds.
\end{apctrap}

\begin{workedexample}[I do: how much crystallises out on cooling]
\textbf{\SI{100}{\gram} of water at \SI{60}{\celsius} is saturated with
\ce{KNO3}, which dissolves to \SI{110}{\gram}/\SI{100}{\gram} at that
temperature. The solution is cooled to \SI{20}{\celsius}, where the
solubility is \SI{31.6}{\gram}/\SI{100}{\gram}. How much solid appears?}

\textbf{What is dissolved at the start.} Saturated at
\SI{60}{\celsius}, so \SI{110}{\gram} of \ce{KNO3} is in solution.

\textbf{What can stay dissolved at the end.} At \SI{20}{\celsius} the
same \SI{100}{\gram} of water can hold only \SI{31.6}{\gram}.

\textbf{The difference has nowhere to go.}
\[ 110 - 31.6 = \mathbf{\SI{78.4}{\gram}}\text{ of \ce{KNO3}
   crystallises out} \]

\textbf{This is how recrystallisation purifies a compound.} Dissolve
the impure solid in hot solvent, cool it, and the bulk of your compound
comes out as crystals while the impurities --- present in far smaller
amounts, and so still below \emph{their} solubility limits --- stay
behind in solution. Filter, and what is on the paper is purer than what
you started with.

\textbf{One caution about the wording.} The question said cooled, not
cooled \emph{slowly and undisturbed}. Cooled carefully, this solution
could instead have become supersaturated and dropped nothing at all
until it was seeded.
\end{workedexample}

\begin{yourturn}[YOUR TURN 8 --- four questions]
Use \ce{KNO3} solubilities of \SI{31.6}{\gram}/\SI{100}{\gram} at
\SI{20}{\celsius} and \SI{169}{\gram}/\SI{100}{\gram} at
\SI{80}{\celsius}.
\begin{enumerate}[leftmargin=*,itemsep=0.5em]
  \item \SI{50}{\gram} of \ce{KNO3} is stirred into \SI{100}{\gram} of
        water at \SI{20}{\celsius}. Is the solution saturated?
        \workspace[1.8cm]
  \item How much \ce{KNO3} crystallises if a solution saturated at
        \SI{80}{\celsius} is cooled to \SI{20}{\celsius}?
        \workspace[1.8cm]
  \item How would you tell a supersaturated solution from a saturated
        one? \workspace[1.8cm]
  \item Why does recrystallisation purify a compound?
        \workspace[1.8cm]
\end{enumerate}
\selfcheck{(a) yes, with \SI{18.4}{\gram} undissolved \quad (b)
\SI{137.4}{\gram} \quad (c) add a seed crystal \quad (d) impurities stay
below their own limits}
\keyonly{\begin{apcanswer}(a) \textbf{Yes.} Only \SI{31.6}{\gram} can
dissolve, so the solution is saturated and $50 - 31.6 =
\mathbf{\SI{18.4}{\gram}}$ remains as undissolved solid at the bottom.
(b) $169 - 31.6 = \mathbf{\SI{137.4}{\gram}}$. (c) \textbf{Add a single
seed crystal.} In a saturated solution it simply sits there; in a
supersaturated one it triggers rapid crystallisation of the whole
excess, and the solution warms as it solidifies. (d) Because the
compound is present in a large amount and its solubility limit is
exceeded on cooling, so it crystallises out --- while the
\textbf{impurities are present in much smaller amounts} and stay below
their own solubility limits, so they remain dissolved and are washed
away with the filtrate.\end{apcanswer}}
\end{yourturn}

% =====================================================================
\section*{\sffamily Ladder 9 \textbullet\ Temperature and solubility}
\zsec{11.3}

Two rules, and they point in opposite directions --- which is why you
must know which kind of solute you have before you answer.

\begin{center}
\begin{tikzpicture}
\begin{axis}[
  width=11.2cm, height=6.4cm,
  xlabel={temperature (\si{\celsius})},
  ylabel={solubility (\si{\gram} per \SI{100}{\gram} water)},
  xmin=0, xmax=100, ymin=0, ymax=270,
  xtick={0,20,40,60,80,100}, ytick={0,50,100,150,200,250},
  grid=major, grid style={apcgray!25},
  legend pos=north west, legend style={font=\footnotesize,
    cells={anchor=west}},
  label style={font=\footnotesize}, tick label style={font=\footnotesize},
]
\addplot[seriesA, thick, mark=*, mark size=1.4pt]
  coordinates {(0,13.3) (20,31.6) (40,63.9) (60,110) (80,169) (100,246)};
\addlegendentry{\ce{KNO3} --- rises steeply}
\addplot[seriesB, thick, mark=square*, mark size=1.4pt]
  coordinates {(0,28.0) (20,34.2) (40,40.1) (60,45.8) (80,51.3)
               (100,56.3)};
\addlegendentry{\ce{KCl} --- rises gently}
\addplot[seriesC, thick, mark=triangle*, mark size=1.8pt]
  coordinates {(0,35.7) (20,35.9) (40,36.4) (60,37.1) (80,38.0)
               (100,39.2)};
\addlegendentry{\ce{NaCl} --- almost flat}
\end{axis}
\end{tikzpicture}
\end{center}

\begin{center}
\begin{tikzpicture}
\begin{axis}[
  width=9.6cm, height=4.8cm,
  xlabel={temperature (\si{\celsius})},
  ylabel={\ce{O2} dissolved (mg/L)},
  xmin=0, xmax=50, ymin=0, ymax=16,
  xtick={0,10,20,30,40,50}, ytick={0,4,8,12,16},
  grid=major, grid style={apcgray!25},
  label style={font=\footnotesize}, tick label style={font=\footnotesize},
]
\addplot[seriesA, thick, mark=*, mark size=1.4pt]
  coordinates {(0,14.6) (10,11.3) (20,9.1) (30,7.6) (40,6.4) (50,5.6)};
\node[font=\footnotesize,anchor=west] at (axis cs:22,12.5)
  {\textbf{gases go the other way}};
\node[font=\footnotesize,anchor=west] at (axis cs:22,10.6)
  {heat drives dissolved gas out};
\end{axis}
\end{tikzpicture}
\end{center}

\begin{workedexample}[I do: reading a solubility curve]
\textbf{Most solids get more soluble as temperature rises} --- but by
wildly differing amounts, and a curve is the only reliable source. Read
the graph above:

\textbf{\ce{KNO3}} climbs from \SI{13.3}{} to
\SI{246}{\gram}/\SI{100}{\gram} --- nearly twentyfold. It is the classic
recrystallisation solute precisely because so much drops out on cooling.

\textbf{\ce{NaCl}} climbs from \SI{35.7}{} to only
\SI{39.2}{\gram}/\SI{100}{\gram} over the same \SI{100}{\celsius}. You
cannot usefully purify salt by recrystallisation from water --- cooling
a saturated solution recovers almost nothing --- which is why sea salt
is obtained by \emph{evaporating} the water instead.

\textbf{Gases always go the other way.} Every gas becomes
\emph{less} soluble as temperature rises, with no exceptions. Warming
gives dissolved gas molecules enough kinetic energy to escape the
liquid, and there is no lattice to break in the first place.

\textbf{Two consequences you can observe.} Bubbles form on the inside of
a pan of water long before it boils --- that is dissolved air coming out
of solution as the water warms. And a power station returning warm
cooling water to a river lowers its dissolved oxygen; a river at
\SI{30}{\celsius} holds only \SI{7.6}{\milli\gram\per\litre} against
\SI{14.6}{} at \SI{0}{\celsius}, which is why \emph{thermal} pollution
kills fish without a single toxic chemical being released.
\end{workedexample}

\begin{yourturn}[YOUR TURN 9 --- four questions]
Read values from the curves above.
\begin{enumerate}[leftmargin=*,itemsep=0.5em]
  \item Roughly how much \ce{KNO3} dissolves in \SI{100}{\gram} of
        water at \SI{40}{\celsius}? \workspace[1.4cm]
  \item Which solid's solubility changes least with temperature?
        \workspace[1.4cm]
  \item What happens to the solubility of a gas as temperature rises?
        \workspace[1.6cm]
  \item Why does warm river water threaten fish even when it is
        chemically clean? \workspace[1.8cm]
\end{enumerate}
\selfcheck{(a) about \SI{64}{\gram} \quad (b) \ce{NaCl} \quad (c) it
falls \quad (d) less dissolved oxygen}
\keyonly{\begin{apcanswer}(a) About \textbf{\SI{64}{\gram}}
(\SI{63.9}{\gram} from the table). (b) \textbf{\ce{NaCl}} --- it rises
only from \SI{35.7}{} to \SI{39.2}{\gram}/\SI{100}{\gram} across the whole
range, which is why salt is recovered by evaporation rather than by
cooling. (c) It \textbf{falls}. Warming gives dissolved gas molecules
enough kinetic energy to escape the liquid, so less stays in solution.
This is true of every gas, without exception. (d) Because warm water
holds \textbf{less dissolved oxygen} --- about \SI{7.6}{} against
\SI{14.6}{\milli\gram\per\litre} between \SI{30}{} and
\SI{0}{\celsius} --- and fish extract oxygen from the water through
their gills. Warming the river suffocates them even though nothing toxic
has been added.\end{apcanswer}}
\end{yourturn}

% =====================================================================
\section*{\sffamily Ladder 10 \textbullet\ Pressure and Henry's law}
\zsec{11.3}

Pressure changes the solubility of \textbf{gases} and, for practical
purposes, nothing else. Solids and liquids are already close-packed and
barely compressible, so squeezing them changes very little.

\[ C = kP \qquad\text{or, comparing two conditions,}\qquad
   \frac{C_{1}}{P_{1}} = \frac{C_{2}}{P_{2}} \]

\begin{center}
\begin{tikzpicture}[font=\sffamily\footnotesize]
  % LOW pressure
  \node[font=\sffamily\small\bfseries] at (1.75,3.85) {low pressure};
  \fill[apcgray!22] (0.7,0.4) rectangle (2.8,1.75);
  \draw[very thick] (0.7,3.2) -- (0.7,0.4) -- (2.8,0.4) -- (2.8,3.2);
  % gas above
  \foreach \p in {(1.0,2.2),(1.7,2.5),(2.4,2.15),(1.35,2.9),(2.15,2.8)} {
    \fill \p circle (0.07);}
  % dissolved
  \foreach \p in {(1.1,0.75),(1.9,1.1),(2.5,0.8)} {
    \fill \p circle (0.07);}
  \draw[-{Stealth[length=2mm]},thick] (1.75,1.9) -- (1.75,1.62);
  \node[align=center] at (1.75,-0.3)
    {few gas particles above\\\textbf{few} dissolved};

  % HIGH pressure
  \begin{scope}[xshift=5.2cm]
    \node[font=\sffamily\small\bfseries] at (1.75,3.85) {high pressure};
    \fill[apcgray!22] (0.7,0.4) rectangle (2.8,1.75);
    \draw[very thick] (0.7,3.2) -- (0.7,0.4) -- (2.8,0.4) -- (2.8,3.2);
    \fill[apcgray!45] (0.7,2.55) rectangle (2.8,2.85);
    \draw[very thick] (0.7,2.55) -- (2.8,2.55);
    \draw[-{Stealth[length=2.5mm]},very thick] (1.75,3.35) -- (1.75,2.95);
    \foreach \p in {(0.95,2.0),(1.35,2.35),(1.8,2.05),(2.25,2.35),
                    (2.6,2.05),(1.15,2.25),(2.0,2.42)} {
      \fill \p circle (0.07);}
    \foreach \p in {(0.95,0.7),(1.35,1.15),(1.75,0.75),(2.2,1.25),
                    (2.6,0.85),(1.15,1.5),(1.95,1.5),(2.5,1.45)} {
      \fill \p circle (0.07);}
    \draw[-{Stealth[length=2mm]},thick] (1.75,1.9) -- (1.75,1.62);
    \node[align=center] at (1.75,-0.3)
      {many gas particles above\\\textbf{many} dissolved};
  \end{scope}

  \node[anchor=west,align=left,font=\sffamily\footnotesize]
    at (9.0,1.9)
    {more collisions with the\\
     surface per second, so more\\
     gas is forced into solution};
\end{tikzpicture}
\end{center}

\begin{workedexample}[I do: the soda can]
\textbf{The solubility of a gas is \SI{1.3e-3}{\Molar} under
\SI{1.0}{\atm} of that gas. What is it under
\SI{4.0}{\atm}?}

Solubility is directly proportional to partial pressure, so
\[ C_{2} = C_{1} \times \frac{P_{2}}{P_{1}}
        = \num{1.3e-3} \times \frac{4.0}{1.0}
        = \mathbf{\SI{5.2e-3}{\Molar}} \]
Four times the pressure, four times the dissolved gas. No other
variable enters.

\textbf{Now the everyday version.} A sealed can of fizzy drink is
bottled under \ce{CO2} at several atmospheres, so a lot of \ce{CO2} is
dissolved. Open it, and the pressure above the liquid drops to the
\ce{CO2} partial pressure of ordinary air --- about
\SI{0.0004}{\atm}. The solubility collapses by a factor of
thousands, the drink is now hugely supersaturated, and the excess
\ce{CO2} escapes as bubbles. Leave it open and it goes flat, because
the process runs to completion.

\textbf{Use partial pressure, not total.} Henry's law depends on the
partial pressure of \emph{that gas} alone. Air is \SI{21}{\percent}
oxygen, so water open to the atmosphere at \SI{1.0}{\atm} total
experiences an oxygen partial pressure of only
\SI{0.21}{\atm} --- which is why dissolved oxygen in a lake is
so much lower than what pure \ce{O2} at \SI{1}{\atm} would give.

\textbf{And why divers care.} At depth the total pressure is far above
\SI{1}{\atm}, so far more nitrogen dissolves in the blood.
Surface too quickly and it comes out of solution as bubbles inside the
body --- decompression sickness, Henry's law applied to a diver.
\end{workedexample}

\begin{yourturn}[YOUR TURN 10 --- four questions]
\begin{enumerate}[leftmargin=*,itemsep=0.5em]
  \item A gas dissolves to \SI{0.020}{\Molar} at \SI{1.0}{\atm}.
        Find its solubility at \SI{3.0}{\atm}.
        \workspace[1.6cm]
  \item A gas dissolves to \SI{0.050}{\Molar} at
        \SI{2.5}{\atm}. Find its solubility at
        \SI{1.0}{\atm}. \workspace[1.6cm]
  \item Why does an open drink go flat? \workspace[1.8cm]
  \item Does raising the pressure increase the solubility of a solid?
        Explain. \workspace[1.8cm]
\end{enumerate}
\selfcheck{(a) \SI{0.060}{\Molar} \quad (b) \SI{0.020}{\Molar} \quad
(c) the \ce{CO2} partial pressure above it collapses \quad (d) no ---
solids are barely compressible}
\keyonly{\begin{apcanswer}(a) $0.020 \times 3.0/1.0 =
\mathbf{\SI{0.060}{\Molar}}$. (b) $0.050 \times 1.0/2.5 =
\mathbf{\SI{0.020}{\Molar}}$. (c) Because opening it drops the \ce{CO2}
partial pressure above the liquid from several atmospheres to the
roughly \SI{0.0004}{\atm} of ordinary air. By Henry's law the
solubility falls in proportion, the drink is left far
\textbf{supersaturated}, and the excess \ce{CO2} leaves as bubbles until
almost none remains. (d) \textbf{No --- essentially not at all.} Solids
and liquids are already close-packed and nearly incompressible, so
pressure barely changes the spacing of their particles. Henry's law
applies to \textbf{gases}, whose volume and therefore whose rate of
collision with the liquid surface pressure genuinely does
change.\end{apcanswer}}
\end{yourturn}

% =====================================================================
\section*{\sffamily Mastery tracker}

Tick a row only if \textbf{all four} YOUR TURN questions were right on
the first attempt. Every row here is assessed --- there is no optional
ladder in this chapter, because the optional material was cut instead.

\renewcommand{\arraystretch}{1.3}
\noindent\begin{tabular}{@{}c p{5.9cm} c p{4.9cm}@{}}
\toprule
\textbf{First try?} & \textbf{Skill} & \textbf{Ladder} &
  \textbf{If not, re-read\ldots}\\
\midrule
$\square$ & molarity & 1 & per litre of \emph{solution} \\
$\square$ & dilution & 2 & moles are conserved \\
$\square$ & mass percent, ppm, ppb & 3 & mass of \emph{solution} \\
$\square$ & mole fraction and molality & 4 & check the denominator \\
$\square$ & energetics of dissolving & 5 & three steps, one exothermic \\
$\square$ & like dissolves like & 6 & can step 3 pay for step 2? \\
$\square$ & chains, heads and tails & 7 & the ratio decides \\
$\square$ & saturated to supersaturated & 8 & solubility is a limit \\
$\square$ & temperature and solubility & 9 & gases go the other way \\
$\square$ & pressure and Henry's law & 10 & partial pressure, gases only \\
\bottomrule
\end{tabular}

\begin{apcnote}[Scoring yourself honestly]
10/10: you have Chapter 11's assessed content. It is a short chapter
only because most of Zumdahl's Chapter 11 was cut from the CED --- what
remains is dense and heavily used elsewhere.

7--9: check \emph{which} ladders missed. Ladders 1--4 are arithmetic and
repair quickly with practice. Ladders 5--7 are conceptual, and a miss
there usually means the three-step cycle has not clicked; that one takes
a re-read rather than more problems.

6 or fewer: start with Ladders 1--3 and do nothing else until they are
automatic. Concentration calculations are not really Chapter 11 content
--- they are the arithmetic of Units 4, 5, 6, 7 and 8 as well. Every
titration, every equilibrium table, every rate law and every cell
potential problem you meet for the rest of the year begins by putting a
concentration into a formula. Time spent here is repaid many times over.
\end{apcnote}

\end{document}
